\chapter{Discussion}
\label{chap:discussion}

\Cref{tab:persistence-synthesis,tab:hypothesis-assessment} show the thesis' central contribution: house-purchase lending growth carries the clearest positive horizon durability, lending spreads provide supporting intermediation evidence, compensation-linked variables are less consistent, and the DAX is sustained in the opposite direction. The discussion interprets this as reduced-form evidence of heterogeneous macro-financial transmission. It is not a claim about household welfare, inequality, redistribution, or structural bank mediation.

\section{Housing-Credit Transmission as the Core Result}

The main empirical result is that ECB policy-news surprises leave the clearest and most persistent trace in house-purchase lending growth. This result is stronger than the evidence for broad credit quantities, liquidity-stock responses, or compensation-linked variables. The housing-credit finding is therefore the center of the thesis rather than a supporting example within a broader asset-price narrative.

This emphasis matters for interpretation. A policy shock that is visible in house-purchase lending growth indicates transmission through a credit-financed, collateral-sensitive sector. That is a macro-financial result. It does not require claiming that the model identifies house-price effects, household access outcomes, or distributional incidence. The evidence is strongest when stated as a channel-specific response: housing-credit variables respond more persistently than the main compensation-linked proxies.

\section{Stocks, Flows, and the Lending Result}

The difference between house-purchase lending growth and pure-new mortgage loans should be interpreted as a stock-flow distinction. Stock-based house-purchase credit growth captures cumulative financing conditions and sustained medium-run accommodation. Flow-based pure-new lending captures immediate origination behavior and is more exposed to short-run volatility, sample length, refinancing composition, and reporting noise.

This distinction reconciles the evidence. The strong annual-growth response shows persistent credit accommodation in the housing-credit stock. The weaker pure-new-loan evidence prevents overstatement of immediate origination effects. Together they support a disciplined mechanism: the housing-credit channel is visible in sustained credit growth, while flow-based originations remain a more volatile robustness surface.

\section{Financial Intermediation Mechanism}

The intermediation evidence makes the housing-credit result economically plausible. Mortgage spreads and NFC spreads remain positive and cumulatively durable, with the mortgage spread providing the most stable lending-condition margin. Broad NFC credit is weaker, which rules out a generic credit-quantity story and points toward sectorally differentiated transmission.

The sequential timing outputs reinforce this reading by ordering lending conditions, credit quantities, housing-credit growth, new mortgage flows, and wage-pressure proxies. The evidence should not be described as structural mediation. The variables respond to the same policy-news shock, and the model does not observe bank-level balance sheets or loan-level decisions. The correct interpretation is timing-consistent financial-channel amplification.

\section{Equity Responses and Information Effects}

The DAX response is negative at medium and long horizons, which is the central financial-market interpretation problem. It should not be ignored, and it should not be forced into a simple asset-appreciation story. The negative response is more consistent with the information-effects literature: ECB easing surprises may partly reveal adverse macroeconomic information during policy communication.

This interpretation is economically plausible. Markets react to both policy accommodation and information about future growth, expected earnings, inflation risks, and risk premia. If the information component is adverse, an easing surprise can coincide with a negative equity response. The DAX therefore acts as a boundary on broad asset-price claims. It supports the view that financial-market responses are mixed and information-sensitive, while the housing-credit and lending-condition evidence carries the main mechanism claim.

\section{Compensation-Linked Responses}

The compensation-linked block is weaker, not absent. The wage tracker excluding one-offs moves at short and intermediate horizons but loses h24 cumulative exclusion. Employment expectations respond earlier and then reverse. The German industry wage-bill proxy supplies a more positive medium-horizon robustness signal, but it is national and sectoral rather than a comprehensive income measure.

The result is therefore comparative. The thesis can state that compensation-linked variables respond less consistently and less persistently than house-purchase lending growth and lending spreads. It cannot state that income gains were absent across households, or that monetary policy caused welfare losses. Those claims require microdata and a different identification design.

\section{Post-COVID and Regime Interpretation}

Regime outputs describe historical dependence rather than replacing the full-sample hierarchy. COVID has too few long-horizon observations, and the tightening regime is short relative to h24. These constraints are important because the thesis is explicitly concerned with persistence. A short regime cannot fully identify a two-year response profile.

Within that scope, the regime evidence is still useful. Tightening shows positive housing-credit and spread accumulation with negative wage-pressure cumulation. COVID shows large short-horizon movements across employment expectations, wage-pressure measures, the DAX, and house-purchase lending growth. These results indicate state sensitivity and support the decision to keep regime evidence descriptive. They enrich the interpretation without overturning the main full-sample claim.

\section{Policy Relevance}

The findings have moderate but real policy relevance. If ECB policy-news shocks transmit persistently through housing credit and lending conditions while compensation-linked proxies respond less consistently, then monetary-policy evaluation should track the duration and sectoral composition of financial transmission. Mortgage conditions, house-purchase credit growth, broad credit stocks, market prices, and compensation proxies should not be collapsed into a single aggregate channel.

These implications remain analytical rather than prescriptive. The thesis does not recommend a particular ECB reaction function, welfare weight, or distributional policy. It shows that monetary transmission has a durable housing-credit and intermediation component, and that this component should be visible in macro-financial analysis.

\section{Limitations}

The thesis' limits define its inferential strength. The design uses high-frequency ECB policy-news surprises, not a fully purified structural monetary-policy innovation. Information effects, event contamination, and regime sensitivity remain part of the identification environment. The first-stage bridge to low-frequency treatments is meaningful in the full sample but sensitive across jackknife and rolling diagnostics, so the thesis does not promote stronger DFR, shadow-rate, or balance-sheet treatment claims.

The aggregation boundary is equally important. Quarterly house prices, compensation per employee, and Bank Lending Survey series are not interpolated into monthly responses. Household-level access metrics, debt-service ratios, borrower heterogeneity, regional housing prices, tenure status, and bank-level balance sheets are outside the empirical design. Statistical limitations also remain: multiple response variables, mixed signs, proxy dependence, and horizon dependence require caution.

Within those boundaries, the conclusion is still substantive. Under a reproducible monthly LP framework, ECB policy-news shocks display more durable evidence in house-purchase lending growth and lending spreads than in compensation-linked proxies and broad productive-credit quantities. That is a disciplined macro-financial transmission result, not a welfare or inequality claim.
